% Generated draft: requires mathematical and source review.
\documentclass[11pt]{article}
\usepackage[margin=25mm]{geometry}
\usepackage{fontspec}
\setmainfont{Latin Modern Roman}
\usepackage{amsmath,amssymb,mathtools,unicode-math}
\setmathfont{Latin Modern Math}
\usepackage{enumitem,needspace,xurl,hyperref}
\setlength{\parindent}{0pt}
\setlength{\parskip}{5pt}
\setlength{\emergencystretch}{3em}
\newcommand{\N}{\mathbb{N}}
\newcommand{\Z}{\mathbb{Z}}
\newcommand{\Q}{\mathbb{Q}}
\newcommand{\R}{\mathbb{R}}
\newcommand{\C}{\mathbb{C}}
\newcommand{\F}{\mathbb{F}}
\newcommand{\A}{\mathbb{A}}
\newcommand{\PP}{\mathbb{P}}
\newcommand{\E}{\mathbb{E}}
\newcommand{\eps}{\varepsilon}
\newcommand{\dd}{\,\mathrm{d}}
\begin{document}
\section*{UnsolvedMath ID: 11000135}
\textbf{Source problem code:} catalogue entry
\subsection{Definitions and mathematical statement}
Fix a surface \(F\) in the source's class. For a pseudo-Anosov (pA) mapping class \(f\), let \(\lambda(f)>1\) be its stretch factor. In the ordinary dilatation-spectrum interpretation, put
\[
 \Sigma(F)=\{\lambda(f):f\text{ is pseudo-Anosov on }F\}.
\]
The problem is to calculate this spectrum. The more modest requests are to determine its least element, or the least gap between distinct elements. To avoid assuming a positive attained gap, one can first ask for
\[
 \Delta(F)=\inf\{|x-y|:x,y\in\Sigma(F),\ x\ne y\},
\]
and whether this infimum is achieved. The source does not supply its preceding definition of \(\Sigma\); if it uses logarithms of dilatations, gaps must instead be computed in that normalization. Existence of pseudo-Anosovs on \(F\), and hence nonemptiness of the spectrum, must also be part of the surface assumptions.

The final request is to characterize number fields arising from dilatations on \(F\), naturally interpreted as
\[
 \{\Q(\lambda(f)):f\text{ pseudo-Anosov on }F\}.
\]
The exact interpretation as generated fields, their abstract isomorphism classes or concrete embeddings is not specified. The source links Problems 4 and 5 and says the recipe in Theorem 4 gives least-element estimates, citing [26, 2]; it also credits McMullen [50] with estimates and examples. The missing recipe and estimates are not reconstructed.
\subsection{Short English statement}
Determine the pseudo-Anosov dilatation spectrum for a fixed surface, or at least its minimum or smallest spacing, and characterize the fields generated by those dilatations.
\subsection{Sources}
\begin{itemize}
\item ``Probing mapping class groups using arcs,'' in \emph{Problems on Mapping Class Groups and Related Topics}, edited by Farb, Problem 5, supplied PDF page 117. \url{https://www.math.uchicago.edu/~farb/papers/mcgbook.pdf}. Recorded access 2026-07-29; extraction pdf-font-structure. The chapter author is not supplied in this record.
\item The source refers to Problem 4, Theorem 4 and [26, 2], and to McMullen [50] for estimates and examples. Full referenced statements and bibliography entries are absent.
\item Local UnsolvedMath registry, \texttt{data/unsolvedmath/problems.json}, permanent ID \(11000135\), exact source code \texttt{catalogue entry}; supplied OPEN-TRIAGE summary reports unverified status.
\end{itemize}
\end{document}
